\documentclass[11pt,a4paper]{article}

\usepackage[utf8]{inputenc}
\usepackage[T1]{fontenc}
% shorthands=off: con 'spanish', las comillas " son un atajo activo, y el texto las usa
\usepackage[spanish,es-noquoting,es-tabla,shorthands=off]{babel}
\usepackage{lmodern}
\usepackage[margin=2.2cm]{geometry}
\usepackage{amsmath,amssymb,amsthm}
\usepackage{booktabs,tabularx,array}
\usepackage{xcolor}
\usepackage{tikz}
\usepackage{pgfplots}
\usepackage{pgfplotstable}
\usepackage{caption}
\usepackage{enumitem}
\usepackage[hidelinks]{hyperref}

\pgfplotsset{compat=1.18}
\usetikzlibrary{positioning,arrows.meta}
\setlist{nosep,leftmargin=*}
\captionsetup{font=small,labelfont=bf}
\decimalpoint
% \code{} no se divide; esto deja estirar los espacios antes que desbordar el margen
\setlength{\emergencystretch}{2.5em}

% Los datos de las figuras se leen directamente de los CSV versionados de los experimentos.
% Se compila desde informe/ (latexmk) o desde la raíz del repo (Overleaf): se prueban ambas rutas.
\IfFileExists{../experiments/results/costs.csv}
  {\newcommand{\resultsdir}{../experiments/results}}
  {\newcommand{\resultsdir}{experiments/results}}
\pgfplotsset{
  discard if not/.style 2 args={
    x filter/.append code={
      \edef\tempa{\thisrow{#1}}\edef\tempb{#2}
      \ifx\tempa\tempb\else\def\pgfmathresult{inf}\fi
    }
  }
}

\newcommand{\code}[1]{\texttt{#1}}
\newcommand{\cause}[1]{\texttt{\small #1}}
\newcommand{\Hash}{\mathsf{H}}
\newtheorem{proposicion}{Proposición}

\title{\textbf{Auditor probabilístico de integridad para almacenamiento en la nube}\\[2pt]
\large Proyecto 7 · Informe final (E2) · Criptografía Aplicada, Uninorte}
\author{Equipo 12: Samuel Camargo Consuegra, Fátima Castro Castro, Juan David Borja Cantillo}
\date{Noviembre de 2026}

\begin{document}
\maketitle

\begin{abstract}
Construimos un auditor que comprueba, sin descargar el archivo, si un servidor conserva sus bloques.
El cliente guarda una huella pequeña y autenticada del archivo (el manifiesto) y envía desafíos con
índices aleatorios; el servidor responde con esos bloques y sus caminos en un árbol de Merkle con
SHA-256 y separación de dominios. Si $d$ de los $n$ bloques están dañados, una auditoría de $q$
índices lo detecta con una probabilidad exacta $P(n,d,q)$ que depende solo de $n$, $d$ y $q$; en 28
combinaciones de daño y tamaño de desafío, con 1000 repeticiones cada una, la tasa observada fue
compatible con esa fórmula en todas. El producto es un auditor por muestreo y no una prueba de
recuperabilidad: un PASS solo dice algo de los bloques consultados.
\end{abstract}

% ------------------------------------------------------------------------------------------------
\section{Qué garantiza el sistema y qué no}
\label{sec:garantia}

El enunciado pide presentar el producto como auditor basado en muestreo y no como prueba formal de
que el archivo es recuperable. Fijamos la garantía en una sola frase, que es la que implementa y
prueba el código:

\begin{quote}
\emph{Si una auditoría termina en PASS, entonces, salvo una colisión de SHA-256, cada uno de los $q$
bloques consultados es idéntico al bloque de esa posición en la versión del archivo comprometida en
el manifiesto. Si el servidor no puede entregar correctos $d$ de los $n$ bloques, una auditoría lo
detecta con probabilidad $P(n,d,q)$.}
\end{quote}

\noindent No se garantiza: (i) que el archivo completo esté disponible o sea recuperable, porque no
hay extractor, códigos de borrado ni autenticadores homomórficos como en los PoR de Juels-Kaliski
\cite{jk07} y Shacham-Waters \cite{sw08}; (ii) nada sobre los bloques no consultados; (iii) que los
datos sigan intactos después de la última auditoría; (iv) disponibilidad ni tiempo de respuesta
(un timeout da INCOMPLETE y no cuenta como evidencia); (v) confidencialidad, pues los bloques viajan
en claro; y (vi) que el servidor guarde los datos él mismo y no los traiga de otra fuente al
responder. La E1 listaba como propiedades \emph{soundness \& extractability}, cifrado de bloques y
sobrecosto de almacenamiento por códigos de redundancia; las quitamos porque el producto no las da
(sección~\ref{sec:desviaciones}).

% ------------------------------------------------------------------------------------------------
\section{Modelo de adversario}
\label{sec:adversario}

\paragraph{Adversario.} El servidor es malicioso y adaptativo y controla el canal. Puede borrar,
modificar, truncar, reordenar o reemplazar bloques y nodos del árbol; conservar versiones anteriores
y responder con ellas; responder con caminos de otro índice, respuestas parciales, duplicadas o de
una auditoría anterior; y ver todos los desafíos pasados antes de decidir qué datos conservar.

\paragraph{Suposiciones.} (S1) SHA-256 es resistente a colisiones (Boneh-Shoup, Def.~8.1
\cite{bs23}). (S2) HMAC-SHA256 es un MAC seguro con la clave local del cliente. (S3) Los índices
salen de un CSPRNG (\code{secrets.SystemRandom}) y el servidor no los predice antes de recibir el
desafío. (S4) La máquina del cliente es confiable: manifiesto, clave y bitácora juegan el papel de
la memoria de solo lectura de Boneh-Shoup §8.1. (S5) La carga inicial ocurre con el servidor honesto.

\paragraph{Fuera del modelo.} Denegación de servicio, un servidor que recupera los datos de otra
fuente al responder, canales laterales, un cliente comprometido y fallos del cliente durante la carga.

% ------------------------------------------------------------------------------------------------
\section{Diseño criptográfico}
\label{sec:diseno}

Todos los parámetros están fijados en un solo módulo (\code{core/params.py}) y en el documento
normativo \code{docs/spec.md}, que el código implementa byte a byte.

\begin{table}[ht]
\centering\small
\begin{tabular}{@{}llll@{}}
\toprule
Prefijo & Función & Entrada a $\Hash$ & Bytes \\
\midrule
\code{0x00} & $\mathsf{leaf}(i,x)$ & $\code{0x00}\,\|\,u_{64}(i)\,\|\,u_{64}(|x|)\,\|\,x$ & $17+|x|$ \\
\code{0x01} & $\mathsf{node}(l,r)$ & $\code{0x01}\,\|\,l\,\|\,r$ & 65 \\
\code{0x02} & $\mathsf{EMPTY\_ROOT}$ & $\code{0x02}\,\|\,\texttt{"empty"}$ & 6 \\
\code{0x03} & $\mathsf{commit}$ & $\code{0x03}\,\|\,\mathit{fid}\,\|\,u_{64}(v)\,\|\,u_{64}(B)\,\|\,u_{64}(S)\,\|\,u_{64}(n)\,\|\,\mathit{root}$ & 81 \\
\bottomrule
\end{tabular}
\caption{Funciones con separación de dominios. $\Hash$ es SHA-256 y $u_{64}$, un entero sin signo de
8 bytes big-endian. Es la única definición de hoja del sistema.}
\label{tab:dominios}
\end{table}

\subsection{Bloques, árbol y compromiso}
El archivo de $S$ bytes se parte en $n=\lceil S/B\rceil$ bloques, con $B\in\{4,8,16,32,64\}$~KiB y
$S\le 500$~MiB. El último bloque es corto y \emph{no se rellena}: su longitud real entra en la hoja.
Cada nivel se forma con $\mathsf{node}$ sobre pares consecutivos; si el nivel es impar, el último
nodo \emph{se promueve} sin volver a aplicar el hash, que es la regla del Ejercicio 8.9 de
Boneh-Shoup (figura~\ref{fig:arbol}). Para $n=0$ la raíz es $\mathsf{EMPTY\_ROOT}$. La huella del
archivo es el compromiso
\[
C=\mathsf{commit}(\mathit{file\_id}, v, B, S, n, \mathit{root}),
\]
que generaliza la raíz final $h(\textit{raíz}\,\|\,s)$ de ese ejercicio: además de $n$ liga el
archivo, la versión y los parámetros. El \code{file\_id} son 16 bytes aleatorios fijos entre
versiones y $v\ge 1$ crece de uno en uno.

\begin{figure}[ht]
\centering
\begin{tikzpicture}[
  every node/.style={draw,rounded corners=2pt,font=\footnotesize,minimum height=5mm,inner sep=2pt},
  lv/.style={draw=none,font=\scriptsize\itshape,text=gray},
  level distance=10mm]
\node (l0) at (0,0) {$\ell_0$}; \node (l1) at (1.4,0) {$\ell_1$};
\node (l2) at (2.8,0) {$\ell_2$}; \node (l3) at (4.2,0) {$\ell_3$};
\node[fill=yellow!25] (l4) at (5.6,0) {$\ell_4$};
\node (a) at (0.7,1.1) {$\mathsf{node}(\ell_0,\ell_1)$};
\node (b) at (3.5,1.1) {$\mathsf{node}(\ell_2,\ell_3)$};
\node[fill=yellow!25] (c) at (5.6,1.1) {$\ell_4$};
\node (d) at (2.1,2.2) {$\mathsf{node}(a,b)$};
\node[fill=yellow!25] (e) at (5.6,2.2) {$\ell_4$};
\node[fill=blue!10] (r) at (3.85,3.3) {$\mathit{root}=\mathsf{node}(d,\ell_4)$};
\draw (l0)--(a); \draw (l1)--(a); \draw (l2)--(b); \draw (l3)--(b);
\draw[dashed] (l4)--(c); \draw (a)--(d); \draw (b)--(d); \draw[dashed] (c)--(e);
\draw (d)--(r); \draw (e)--(r);
\node[lv] at (8.2,1.1) {$\ell_4$ se promueve};
\node[lv] at (8.2,2.2) {y otra vez};
\end{tikzpicture}
\caption{Árbol para $n=5$. El nodo sin hermano sube tal cual (línea discontinua). El camino de
$\ell_4$ tiene un solo hash ($d$), y su forma la calcula el verificador a partir de $(4,5)$.}
\label{fig:arbol}
\end{figure}

\subsection{Pruebas de inclusión}
El camino de la hoja $i$ es la lista de hashes hermanos desde la hoja hacia la raíz. \emph{La
dirección de cada hermano se deriva de $(i,n)$ y nunca del servidor}: la función
$\mathsf{path\_shape}(i,n)$ recorre los niveles y, en cada uno, devuelve LEFT si el índice actual es
impar, RIGHT si tiene hermano a la derecha, o nada si el nodo se promueve. Verificar un bloque
consiste en (1) comprobar $|x|=L(i)$, donde $L(i)=B$ salvo el último bloque, (2) comprobar que el
camino tiene la longitud que da $\mathsf{path\_shape}$, (3) recomputar la raíz desde
$\mathsf{leaf}(i,x)$ y (4) compararla en tiempo constante con la del manifiesto. Cada paso tiene su
causa de rechazo (\cause{BLOCK\_LENGTH}, \cause{PATH\_LENGTH}, \cause{ROOT\_MISMATCH}).

\subsection{Por qué la codificación es inequívoca}
La raíz funciona como un \emph{commitment} del archivo, y el binding por resistencia a colisiones
exige una codificación inequívoca de lo que se compromete (\cite{vp}, pág.~49--52). Lo
justificamos así: hojas y nodos internos se distinguen por el primer byte, como en Certificate Transparency \cite{rfc9162}; posición y longitud entran
con ancho fijo, de modo que el mismo contenido en otra posición, o el mismo bloque con ceros
añadidos, dan hojas distintas; la forma del árbol depende solo de $n$, que está ligado en $C$, y el
verificador nunca toma del servidor qué hash va a cada lado; y $\mathsf{commit}$ liga
$\mathit{file\_id}$, $v$, $B$, $S$ y $n$ con anchos fijos.

\begin{proposicion}
Si un servidor produce, para un índice $i$ consultado, un par $(x',\pi')$ con $x'\ne x_i$ que pasa
la verificación contra la raíz del manifiesto, entonces se puede extraer una colisión de SHA-256.
\end{proposicion}
\noindent\emph{Idea de la prueba.} Es el argumento de Boneh-Shoup, Thm~8.8, adaptado al árbol del
Ej.~8.9. Como la forma del camino depende solo de $(i,n)$, la recomputación honesta y la falsa
recorren las mismas posiciones. Como ambas terminan en la misma raíz y empiezan en entradas
distintas ($\code{0x00}\|u_{64}(i)\|u_{64}(|x'|)\|x' \ne \code{0x00}\|u_{64}(i)\|u_{64}(|x_i|)\|x_i$),
existe un primer nivel en que las salidas de $\Hash$ coinciden y las entradas difieren, y eso es una
colisión. Los prefijos impiden que una hoja se haga pasar por nodo interno o viceversa. \hfill$\square$

\subsection{Manifiesto autenticado}
El manifiesto guarda \code{file\_id}, versión, nombre, $S$, $B$, $n$, raíz, $C$ y fecha, serializados
como JSON canónico (claves ordenadas, sin espacios, solo ASCII), y lleva un HMAC \cite{rfc2104}
\[
\mathit{mac}=\mathrm{HMAC\text{-}SHA256}\bigl(K,\ \texttt{"merkle-audit/manifest/1"}\,\|\,\code{0x00}\,\|\,\mathsf{canon}(m)\bigr)
\]
con una clave $K$ de 32 bytes que se crea en el primer uso con permisos \code{0600} y nunca se
registra. Antes de cada auditoría se valida en orden: forma, MAC (con
\code{hmac.compare\_digest}), coherencia de parámetros y compromiso. Si algo falla, la auditoría es
INVALID y \emph{no se envía el desafío}. El cliente no conserva el archivo después de la carga: el
manifiesto es todo lo que tiene.

\subsection{Frescura y ligadura al contexto}
Cada auditoría usa un \code{audit\_id} nuevo de 128 bits y $q$ índices distintos, elegidos sin
reemplazo con el generador criptográfico del sistema (\code{secrets.SystemRandom}). La respuesta debe repetir \code{audit\_id}, \code{file\_id} y
versión, y traer exactamente los índices pedidos en el mismo orden; si no, es INVALID. Con eso una
respuesta de otra auditoría, de otro archivo o de otra versión no se acepta
(\cite{vp}, págs.~43 y 142). Un rollback que dice la versión vieja se rechaza por
\cause{VERSION\_MISMATCH}; uno que miente y dice la versión pedida falla por \cause{ROOT\_MISMATCH},
porque sus bloques no coinciden con la raíz comprometida en $C$.

% ------------------------------------------------------------------------------------------------
\section{Protocolo}
\label{sec:protocolo}

El transporte es HTTP local (FastAPI y httpx). La versión del protocolo va en el cuerpo
(\code{"protocol": "merkle-audit/1"}); los modelos de mensaje rechazan campos de más o de menos y
conversiones implícitas de tipos. Hashes e identificadores van en hex en minúsculas de longitud
exacta y los bloques en base64 estricto.

\paragraph{Carga.} \code{PUT /files/\{file\_id\}/versions/\{v\}} con el archivo crudo. El servidor
guarda bloques, niveles del árbol y una copia íntegra (para que el inyector pueda restaurar), y
responde con un recibo que trae $n$, raíz y $C$. El cliente calcula su propio árbol mientras envía y
solo crea el manifiesto si el recibo coincide. Las versiones son inmutables y consecutivas.

\paragraph{Auditoría.} \code{POST /audit} con el desafío
$(\mathit{audit\_id},\mathit{file\_id},v,C,\mathit{indices})$; la respuesta trae una prueba
$(i, x_i, \pi_i)$ por índice, o \code{block: null} si el servidor declara que no tiene ese bloque.
El servidor lee los caminos de los niveles guardados en la carga, no los recalcula, de modo que un
bloque alterado en disco aparece como \cause{ROOT\_MISMATCH} y no como un camino coherente con la
alteración.

\paragraph{Estados.} La evaluación sigue un orden fijo: validar el manifiesto; enviar y esperar
(timeout de 60~s configurable); validar estructura y ligadura de la respuesta (el primer fallo da
INVALID y no se verifica ninguna prueba); y verificar \emph{todas} las pruebas sin parar en la
primera que falla, para reportar cada índice dañado. El código HTTP acompaña pero no decide: el
veredicto sale del cuerpo. La tabla~\ref{tab:estados} resume los cuatro estados y las 16 causas, que
forman una lista cerrada.

\begin{table}[ht]
\centering\small
\begin{tabularx}{\linewidth}{@{}l>{\raggedright\arraybackslash}Xc@{}}
\toprule
Estado & Causas & ¿Detección? \\
\midrule
PASS & --- (con $n=0$ se marca \emph{trivial}) & No \\
CORRUPT & \cause{BLOCK\_MISSING}, \cause{BLOCK\_LENGTH}, \cause{PATH\_LENGTH}, \cause{ROOT\_MISMATCH},
  \cause{FILE\_MISSING} & Sí \\
INVALID & \cause{MANIFEST\_FORMAT}, \cause{MANIFEST\_MAC}, \cause{MANIFEST\_COMMITMENT},
  \cause{MALFORMED}, \cause{AUDIT\_ID\_MISMATCH}, \cause{FILE\_ID\_MISMATCH},
  \cause{VERSION\_MISMATCH}, \cause{INDEX\_SET\_MISMATCH}, \cause{SERVER\_ERROR} & Sí \\
INCOMPLETE & \cause{TIMEOUT}, \cause{CONNECTION} & No \\
\bottomrule
\end{tabularx}
\caption{Estados de una auditoría y lista cerrada de causas. INVALID cuenta como detección porque el
servidor no cumplió el protocolo; INCOMPLETE no aporta evidencia en ningún sentido.}
\label{tab:estados}
\end{table}

% ------------------------------------------------------------------------------------------------
\section{Modelo probabilístico}
\label{sec:prob}

Si $d$ de los $n$ bloques están dañados y se consultan $q$ índices uniformes sin reemplazo, la
auditoría no ve ninguno dañado con probabilidad $\binom{n-d}{q}/\binom{n}{q}$, así que
\[
P(n,d,q)=1-\frac{\binom{n-d}{q}}{\binom{n}{q}}=1-\prod_{j=0}^{q-1}\frac{n-d-j}{n-j},
\]
con $P=0$ si $d=0$ o $q=0$, y $P=1$ si $q>n-d$. La fórmula es exacta en este modelo porque el
servidor fija qué bloques pierde antes de ver el desafío y no puede predecirlo (S3). El código usa
la forma de producto en coma flotante; los tests la comparan con \code{math.comb} y
\code{fractions.Fraction} para $n\le 2000$. La expresión $1-(1-\varepsilon)^q$ de la E1 corresponde al
muestreo \emph{con} reemplazo y solo la usamos como comparación.

\paragraph{Elección de $q$.} $q_{\min}(n,d,1-\delta)$ es el menor $q$ con $P\ge 1-\delta$; como $P$
crece con $q$, se halla por bisección. Por defecto el objetivo es detectar $\varepsilon=1$\,\% con
99~\%. Para $n=10\,000$ eso da $q=448$ ($P=99{,}002$\,\%), frente a 459 con la aproximación. El valor
casi no depende de $n$: 368 para $n=1000$, 455 para $n=25\,600$ (100~MiB con bloques de 4~KiB) y
458 para el máximo, $n=128\,000$. Con $Q_{\max}=1024$ todos caben en una respuesta.

\paragraph{Repetición y límites.} Con $k$ auditorías independientes y daño estático, la probabilidad
de no detectar es $(1-P)^k$. Esto no prueba disponibilidad: el daño puede ocurrir después de la
última auditoría; un daño de un solo bloque se ve por auditoría con probabilidad $q/n$ (4,5~\% con
$q=448$ y $n=10\,000$); y los bloques nunca consultados no tienen ninguna garantía. La demostración
reproducible muestra justamente este caso: un bloque modificado en un archivo de 2048 bloques suele
pasar una auditoría con $q=300$.

% ------------------------------------------------------------------------------------------------
\section{Implementación}
\label{sec:impl}

El código (Python 3.12, dependencias fijadas con \code{uv.lock}) se organiza en capas con una regla
de dependencias estricta, $\code{core}\leftarrow\code{protocol}\leftarrow\{\code{client},\code{server}\}\leftarrow\code{ui}$,
como pide el enunciado:

\begin{itemize}
\item \code{core}: primitivas puras sin red ni disco: codificación, árbol, manifiesto y probabilidad.
\item \code{protocol}: mensajes (pydantic, \code{extra="forbid"}, modo estricto), generación de
  desafíos y evaluación de respuestas.
\item \code{server}: API FastAPI y almacén SQLite (bloques, niveles, copia íntegra y fallos, en modo
  WAL), más el inyector de fallos \code{merkle-audit-fault}, un CLI sin endpoint de red que escribe
  en el almacén del servidor y nunca en el cliente. Implementa borrado, modificación (bit invertido o
  truncado), reemplazo, camino cruzado, rollback declarado o mentido, daño de una fracción $\varepsilon$,
  pérdida del archivo y restauración; cada aplicación queda registrada con su semilla.
\item \code{client}: CLI \code{merkle-audit} (\code{upload}, \code{audit}, \code{reverify},
  \code{log}, \code{schedule}), transporte httpx, almacén de manifiestos y clave, programador de
  auditorías periódicas y la bitácora.
\item \code{ui}: panel Streamlit con las curvas de probabilidad y de costo en función de $q$, el
  historial con la evidencia de cada auditoría y los programas.
\end{itemize}

\paragraph{Bitácora.} Cada auditoría queda en una tabla SQLite \emph{solo de inserción} (dos
triggers impiden \code{UPDATE} y \code{DELETE}) con identidad, parámetros, la lista completa de
índices, el estado y las causas, la evidencia por índice (hoja y raíz recomputadas en hex), el
SHA-256 del cuerpo de la respuesta, latencia, bytes y tiempo de verificación, y $P(n,d,q)$ para
$\varepsilon\in\{1,5,10,25\}$\,\%. No se guardan bloques ni la clave. Con \code{-{}-keep-responses}
se guarda la respuesta cruda y \code{reverify} repite la verificación; esta opción es solo para datos
sintéticos.

\paragraph{Aleatoriedad en los experimentos.} Solo el código de \code{experiments/} puede usar como fuente
de índices un generador con semilla (\code{random.Random}), para que los resultados se puedan repetir. El
CLI y el panel no exponen esa opción.

% ------------------------------------------------------------------------------------------------
\section{Pruebas}
\label{sec:pruebas}

La suite tiene 295 tests (\code{make check} corre lint y tests; uno se omite en Windows porque
comprueba permisos POSIX): 153 unitarios, 68 de ataques y 74 de integración. Cada regla de rechazo
tiene su test en \code{tests/adversarial/}; los de integración prueban cliente y servidor reales de
punta a punta y la demostración completa. La tabla~\ref{tab:pruebas} resume la trazabilidad.

\begin{table}[ht]
\centering\small
\begin{tabularx}{\linewidth}{@{}>{\raggedright\arraybackslash}X>{\raggedright\arraybackslash}X@{}}
\toprule
Requisito o causa & Test \\
\midrule
Vectores de la spec (vacío, 1 byte, 3 bloques) & \code{unit/test\_vectors.py} \\
Forma del camino: verifica para todo $i<n\le 300$ (hypothesis) & \code{unit/test\_merkle.py} \\
$P$ exacta y $q_{\min}$ & \code{unit/test\_probability.py} \\
Archivos intactos, vacíos y de tamaño no múltiplo de $B$ & \code{integration/test\_server\_api.py}, \code{test\_client.py} \\
\cause{MANIFEST\_FORMAT}, \cause{MANIFEST\_MAC}, \cause{MANIFEST\_COMMITMENT} & \code{adversarial/test\_manifest\_tampering.py} \\
\cause{MALFORMED} (19 casos de esquema, hex y base64) & \code{adversarial/test\_malformed.py} \\
\cause{AUDIT\_ID\_MISMATCH} (replay), \cause{FILE\_ID\_MISMATCH} & \code{test\_replay.py}, \code{test\_binding.py} \\
\cause{VERSION\_MISMATCH}, rollback mentido & \code{adversarial/test\_rollback.py} \\
\cause{INDEX\_SET\_MISMATCH}: falta, sobra, duplicado, ajeno, orden & \code{adversarial/test\_index\_set.py} \\
\cause{BLOCK\_MISSING}, \cause{BLOCK\_LENGTH}, \cause{ROOT\_MISMATCH} & \code{test\_delete.py}, \code{test\_modify.py}, \code{test\_replace.py} \\
Camino cruzado (\cause{ROOT\_MISMATCH} / \cause{PATH\_LENGTH}) & \code{adversarial/test\_cross\_path.py} \\
\cause{SERVER\_ERROR}, \cause{FILE\_MISSING} & \code{adversarial/test\_server\_errors.py} \\
\cause{TIMEOUT}, \cause{CONNECTION}; manifiesto alterado sin enviar el desafío & \code{integration/test\_client.py} \\
Bitácora solo de inserción, programador, \code{reverify} & \code{integration/test\_client.py} \\
\bottomrule
\end{tabularx}
\caption{Trazabilidad de requisitos y causas a pruebas automatizadas (rutas bajo \code{tests/}).}
\label{tab:pruebas}
\end{table}

Además, \code{make demo} recorre con los CLI reales los casos del enunciado (intacto, vacío, tamaño
impar, borrado, modificación, reemplazo, camino cruzado, rollback declarado y mentido, pérdida del
archivo, manifiesto alterado, reverificación y servidor caído), declara el estado que espera en cada
paso y termina con error si alguno no coincide. En la última ejecución, los 20 pasos salieron como se
esperaba.

% ------------------------------------------------------------------------------------------------
\section{Experimentos}
\label{sec:exp}

Los scripts de \code{experiments/} son reproducibles con semilla (2026 por defecto) y escriben CSV
que se versionan; las figuras y tablas de esta sección se generan directamente de esos CSV al compilar.

\subsection{Detección: empírica frente a teórica}
\paragraph{Diseño.} Archivo sintético de $n=1000$ bloques de 4~KiB. Para cada
$\varepsilon\in\{1,5,10,25\}$\,\% se dañan $d=\lceil\varepsilon n\rceil$ bloques (un bit invertido)
y para cada $q\in\{5,10,25,50,100,200,300\}$ se hacen $N=1000$ auditorías con índices nuevos. El
servidor simulado conserva el árbol original y cada respuesta pasa por la misma evaluación del
cliente; cuenta como detección toda auditoría que no termina en PASS. Una combinación se acepta si
$P(n,d,q)$ cae dentro del intervalo de Wilson \cite{wilson} al 99~\% de la tasa observada $\hat p$.

\paragraph{Resultado.} Las 28 combinaciones quedan dentro de su intervalo
(figura~\ref{fig:deteccion}). Con 28 pruebas al 99~\% se esperaban en promedio 0,3 fuera, así que el
resultado es el esperable. Un ejemplo: con $\varepsilon=1$\,\% y $q=100$, $\hat p=0{,}665$ frente a
$P=0{,}653$, con intervalo $[0{,}626;\ 0{,}702]$. La E1 usaba una tolerancia fija de $\pm1{,}5$\,\%,
que en ese punto es apenas una desviación estándar y habría rechazado cerca de un tercio de las
combinaciones correctas.

\begin{figure}[ht]
\centering
\begin{tikzpicture}
\begin{axis}[
  width=0.86\linewidth, height=6.2cm,
  xmode=log, log basis x=10,
  xtick={5,10,25,50,100,200,300}, xticklabels={5,10,25,50,100,200,300},
  xlabel={$q$ (bloques consultados), $n=1000$}, ylabel={probabilidad de detección},
  ymin=0, ymax=1.03, grid=major, grid style={gray!20},
  legend pos=south east, legend cell align=left, legend style={font=\scriptsize},
  cycle list={{blue!70!black},{teal!70!black},{orange!85!black},{red!75!black}},
]
\foreach \e in {0.01, 0.05, 0.1, 0.25} {
  \edef\tmp{
    \noexpand\addplot+[no marks, thick, discard if not={eps}{\e}]
      table[x=q, y=exact, col sep=comma] {\resultsdir/detection.csv};
  }\tmp
}
\legend{$\varepsilon=1$\,\%, $\varepsilon=5$\,\%, $\varepsilon=10$\,\%, $\varepsilon=25$\,\%}
\foreach \e in {0.01, 0.05, 0.1, 0.25} {
  \edef\tmp{
    \noexpand\addplot+[only marks, mark=*, mark size=1.4pt, black, forget plot,
      discard if not={eps}{\e},
      error bars/.cd, y dir=both, y explicit]
      table[x=q, y=empirical, col sep=comma,
        y error plus expr=\noexpand\thisrow{wilson99_high}-\noexpand\thisrow{empirical},
        y error minus expr=\noexpand\thisrow{empirical}-\noexpand\thisrow{wilson99_low}]
      {\resultsdir/detection.csv};
  }\tmp
}
\end{axis}
\end{tikzpicture}
\caption{Probabilidad exacta $P(n,d,q)$ (líneas) y tasa observada en 1000 auditorías (puntos) con su
intervalo de Wilson al 99~\%. Datos: \code{experiments/results/detection.csv}.}
\label{fig:deteccion}
\end{figure}

\subsection{Costos}
\paragraph{Diseño.} Para archivos de 1, 10 y 100~MiB y bloques de 4, 16 y 64~KiB se construye el
árbol, se elige $q_{\min}$ para detectar $\varepsilon=1$\,\% con 99~\% y se mide una auditoría honesta
completa: bytes del desafío y de la respuesta (JSON tal como viaja), tiempo de verificación y
cociente frente a descargar el archivo. Los tiempos son la mediana de 5 repeticiones en una máquina
de desarrollo y solo sirven para comparar entre sí.

\begin{table}[ht]
\centering\footnotesize\setlength{\tabcolsep}{4pt}
\pgfplotstabletypeset[
  col sep=comma,
  columns={file_mib,block_size,n,tree_levels,build_ms,q,avg_path_hashes,audit_bytes,bandwidth_ratio,verify_ms},
  columns/file_mib/.style={column name={MiB}, fixed, precision=0},
  columns/block_size/.style={column name={$B$ (KiB)}, preproc/expr={##1/1024}, fixed, precision=0},
  columns/n/.style={column name={$n$}, int detect, 1000 sep={\,}},
  columns/tree_levels/.style={column name={niv.}},
  columns/build_ms/.style={column name={árbol ms}, fixed, precision=0},
  columns/q/.style={column name={$q$}},
  columns/avg_path_hashes/.style={column name={hashes}, fixed, precision=1, zerofill},
  columns/audit_bytes/.style={column name={MB},
    preproc/expr={##1/1000000}, fixed, precision=2, zerofill},
  columns/bandwidth_ratio/.style={column name={cociente}, fixed, precision=2, zerofill},
  columns/verify_ms/.style={column name={verif. ms}, fixed, precision=0},
  every head row/.style={before row=\toprule, after row=\midrule},
  every last row/.style={after row=\bottomrule},
  every nth row={3}{before row=\midrule},
]{\resultsdir/costs.csv}
\caption{Costo de una auditoría con $q_{\min}$ para $\varepsilon=1$\,\% y 99~\%. «niv.»: niveles del
árbol; «hashes»: longitud media del camino; «MB»: desafío más respuesta; «cociente»: MB de la
auditoría entre el tamaño del archivo; tiempos en ms. Datos: \code{experiments/results/costs.csv}.}
\label{tab:costos}
\end{table}

\paragraph{Resultado.} La tabla~\ref{tab:costos} muestra tres hechos. Primero, como $q_{\min}$ casi
no crece con $n$, el costo de auditar es casi constante y la ventaja frente a descargar crece con el
archivo: con 100~MiB y bloques de 4~KiB una auditoría mueve unos 3~MB, el 2,8~\% del archivo.
Segundo, en archivos pequeños el muestreo no ahorra nada: con 1~MiB hay que consultar 201 de 256
bloques y, por el 33~\% de base64 y los caminos, la auditoría cuesta más que el archivo (cociente
1,16). Para esos tamaños, descargar y comparar con la raíz es más barato y da certeza. Tercero, el
tamaño de bloque domina el ancho de banda: con 64~KiB el camino es más corto pero cada bloque
consultado pesa 16 veces más, y en 100~MiB la auditoría pasa del 2,8~\% al 34~\% del archivo. Por eso
dejamos 4~KiB por defecto. Construir el árbol cuesta del orden de 3~ms por MiB y verificar una
auditoría, menos de un segundo en todos los casos medidos.

% ------------------------------------------------------------------------------------------------
\section{Desviaciones respecto a la propuesta (E1)}
\label{sec:desviaciones}

El enunciado permite corregir el diseño inicial siempre que se expliquen las desviaciones. Cada
cambio se justifica con las correcciones dadas por el docente y con el texto; la tabla~\ref{tab:desv} resume los
que alteran lo propuesto en la E1, incluidos los que señaló la evaluación.

\begin{table}[ht]
\centering\small
\begin{tabularx}{\linewidth}{@{}>{\raggedright\arraybackslash}p{0.27\linewidth}>{\raggedright\arraybackslash}X@{}}
\toprule
E1 & E2 y motivo \\
\midrule
Propiedades: \emph{soundness \& extractability}, cifrado de bloques, sobrecosto por redundancia &
Se quitan. El producto es un auditor por muestreo; los PoR quedan como estado del arte
(sección~\ref{sec:garantia}). \\
Dos definiciones de hoja: $\Hash(\code{0x00}\|i\|\,|x_i|\,\|x_i)$ y $h(x_i)$ &
Una sola $\mathsf{leaf}(i,x)$ con anchos fijos $u_{64}$, en cliente, servidor, tests e informe. \\
Nivel impar «por ejemplo, promover»; último bloque con ceros &
Promoción sin rehash (Ej.~8.9) y último bloque corto sin relleno: con ceros, un archivo y el mismo
con ceros añadidos tenían la misma raíz. \\
Archivo vacío: «error o hash nulo» & $\mathsf{EMPTY\_ROOT}$ con dominio propio y auditoría con $q=0$
que pasa de forma trivial y lo dice. \\
Raíz sola como huella & Compromiso $C$ con prefijo \code{0x03} que liga archivo, versión,
parámetros y $n$. \\
Manifiesto alterado detectado porque «las pruebas no cuadran» & HMAC-SHA256 sobre JSON canónico:
se rechaza antes de enviar el desafío. \\
Camino sin dirección & La dirección sale de $(i,n)$, nunca del servidor. \\
Índices con «secrets / random» & Solo \code{secrets.SystemRandom}; \code{random} con semilla, solo
en experimentos. \\
Respuestas en «lotes 1..K» sin formato & Una respuesta por desafío con $Q_{\max}=1024$, y reglas
cerradas para falta, sobra, duplicado, ajeno y orden. \\
Comparación con $1-(1-\varepsilon)^q$ y tolerancia $\pm1{,}5$\,\% & Probabilidad exacta
hipergeométrica, $q_{\min}$ por bisección y criterio de Wilson al 99~\%. \\
Ejemplo con $n=10^6$ bloques & Imposible con 500~MB y 4~KiB: el máximo es $n=128\,000$, con caminos de a lo sumo 17 hashes. \\
\bottomrule
\end{tabularx}
\caption{Desviaciones entre la propuesta (E1) y la implementación (E2).}
\label{tab:desv}
\end{table}

Además decidimos durante la implementación, sin cambiar nada de la E1: SHA-256 como único hash, sin
negociación; rutas HTTP sin prefijo de versión (la versión va en el cuerpo); códigos HTTP fijos con
el cuerpo como única fuente de verdad; versiones consecutivas; y caminos leídos de los niveles
guardados en la carga, no recalculados.

% ------------------------------------------------------------------------------------------------
\section{Limitaciones y trabajo futuro}
\label{sec:limites}

\begin{itemize}
\item \textbf{No es un PoR.} Sin códigos de borrado ni extractor, un servidor que pierde un solo
  bloque pasa la mayoría de las auditorías. Un PoR al estilo Juels-Kaliski o Shacham-Waters codifica
  el archivo para que perder poco sea irrecuperable y perder mucho sea detectable.
\item \textbf{Ancho de banda.} Cada bloque consultado viaja entero. Las multipruebas
  (Boneh-Shoup, Thm~8.7) reducirían los hashes de los caminos, pero el costo lo dominan los bloques.
  Los autenticadores homomórficos de Shacham-Waters responden con un agregado de tamaño constante.
\item \textbf{Archivos pequeños.} Por debajo de unos pocos MiB, auditar cuesta tanto como descargar
  (sección~\ref{sec:exp}). El producto no lo advierte todavía.
\item \textbf{Confianza en el cliente.} La seguridad descansa en que manifiesto, clave y bitácora no
  se alteren (S4). La bitácora es solo de inserción a nivel SQLite, pero no está encadenada con
  hashes: quien controla la máquina del cliente puede reescribirla.
\item \textbf{Sin actualización dinámica.} Cada cambio del archivo es una versión nueva completa.
\end{itemize}

% ------------------------------------------------------------------------------------------------
\section{Reproducibilidad}
\label{sec:repro}

Desde cero, con \code{uv} instalado: \code{uv sync} crea el entorno con las dependencias fijadas;
\code{make check} corre lint y tests; \code{make demo} ejecuta la demostración de punta a punta;
\code{make experiments} regenera los CSV y sus gráficas (unos 10 minutos con $N=1000$);
\code{make server} y \code{make ui} levantan el servidor y el panel. El README tiene además un guion
manual paso a paso con el CLI y el inyector de fallos. Todos los datos son sintéticos y salen de una
semilla; ningún secreto se registra.

% ------------------------------------------------------------------------------------------------
\begin{thebibliography}{9}\small
\bibitem{bs23} D.~Boneh y V.~Shoup. \emph{A Graduate Course in Applied Cryptography}, versión 0.6,
  enero de 2023. §8.1, §8.9 (Def.~8.3, Thm~8.7, Thm~8.8) y Ejercicio 8.9.
\bibitem{jk07} A.~Juels y B.~S.~Kaliski. PORs: Proofs of Retrievability for Large Files. En
  \emph{ACM CCS}, 2007.
\bibitem{sw08} H.~Shacham y B.~Waters. Compact Proofs of Retrievability. En \emph{ASIACRYPT},
  LNCS 5350, 2008.
\bibitem{vp} R.~L.~Villanueva-Polanco. \emph{Cryptographic Protocols}, material del curso
  Criptografía Aplicada, Universidad del Norte, 2026. Págs.~7--9, 43, 49--52, 113, 142, 248--249.
\bibitem{rfc9162} B.~Laurie, E.~Messeri y R.~Stradling. Certificate Transparency Version 2.0.
  RFC 9162, 2021.
\bibitem{rfc2104} H.~Krawczyk, M.~Bellare y R.~Canetti. HMAC: Keyed-Hashing for Message
  Authentication. RFC 2104, 1997.
\bibitem{wilson} E.~B.~Wilson. Probable Inference, the Law of Succession, and Statistical
  Inference. \emph{Journal of the American Statistical Association}, 22(158):209--212, 1927.
\end{thebibliography}

\end{document}
